%! Author = sriadityavedantam
%! Date = 4/8/26

\section{Behavior Under Maps of Differential 1-Forms}

Let $X,Y$ be quasi-projective varieties and let $\phi: X \to Y$ be a regular map.
There exists a pullback map on $1$-forms
\[
    \phi^\upstar : \Omega[Y] \to \Omega[X].
\]

\subsection*{Example}

Let $\phi : \aa^1_t \to \aa^1_u$ be defined by $t \mapsto u = t^k$.
Then
\[
    \phi^\upstar : \Omega[\aa^1_u] \to \Omega[\aa^1_t]
\]
is given by
\[
    f(u)\,du \longmapsto f(t^k)\,d(t^k),
    \quad
    d(t^k) = k t^{k-1} dt.
\]

\begin{theorem*}{
    If $X,Y$ are smooth projective varieties and birational, then
    \[
        \Omega[X] \cong \Omega[Y].
    \]
}
\end{theorem*}

If $X$ is birational to $Y$, there exists a rational map $\phi: X \dashrightarrow Y$.
Then $\phi$ is regular on some open subset $U \subseteq X$.
For $\omega \in \Omega[Y]$, the pullback
\[
    (\phi|_U)^\upstar \omega
\]
is a regular $1$-form on $U$, hence a rational $1$-form on $X$.

\begin{theorem*}{
    If $\phi$ is birational, then $(\phi|_U)^\upstar \omega$ extends to a regular $1$-form on $X$.
}
\end{theorem*}

Thus $\phi^\upstar : \Omega[Y] \to \Omega[X]$ is well-defined even when $\phi$ is only rational, provided $X,Y$ are smooth projective varieties.
Consequently,
\[
    \dim \Omega[X] = \dim \Omega[Y],
\]
so these dimensions (the Hodge numbers $h^{1,0}$) are birational invariants.

\subsection*{Canonical Divisor}

Let $\omega,\omega' \in \Omega(X)$ be nonzero rational $1$-forms.
Since $\Omega(X)$ is a $1$-dimensional vector space over $\c(X)$ (when $X$ is a curve), we have
\[
    \omega' = g\omega
\]
for some $g \in \c(X)$.
Thus
\[
    \div(\omega') = \div(g) + \div(\omega).
\]
Since $\div(g)$ is principal, we obtain
\[
    \div(\omega) \sim \div(\omega').
\]
Hence $\div(\omega)$ defines a well-defined class in $Cl(X)$, called the \textit{canonical class} $K_X$.

\subsection*{Relation with $\Ell(K_X)$}

For a divisor $D$, recall
\[
    \Ell(D) = \{0\} \cup \{f \in \c(X) : \div(f) + D \geq 0\}.
\]

Then
\[
    \Ell(K_X)
    = \{0\} \cup \{f \in \c(X) : \div(f) + K_X \geq 0\}.
\]

Since $K_X = \div(\omega)$,
\[
    \div(f) + \div(\omega) \geq 0
    \quad \Longleftrightarrow \quad
    \div(f\omega) \geq 0.
\]
Thus $f\omega$ is a regular $1$-form, so
\[
    \Ell(K_X) \cong \Omega[X].
\]

\subsection*{Example: $X = \pp^1$}

Let $\pp^1 = \aa^1_t \cup \aa^1_u$ with $t = \frac{1}{u}$.
Take $\omega = dt$, which is a rational $1$-form on $\pp^1$.
On $\aa^1_t$, we have
\[
    \div(\omega)|_{\aa^1_t} = 0.
\]

On $\aa^1_u$, since
\[
    t = \frac{1}{u},
    \quad
    dt = -\frac{1}{u^2} du,
\]
we obtain
\[
    \div(\omega)|_{\aa^1_u} = \div(-u^{-2}) = -2\{u=0\}.
\]

Hence
\[
    \deg(K_X) = -2.
\]

\subsection*{Problem}

Show that the rational differential form $\frac{dx}{y}$ is regular on the affine curve
\[
    X = \{x^2 + y^2 = 1\}.
\]

If $y \neq 0$, then $\frac{dx}{y}$ is clearly regular.
If $y = 0$, this occurs at two points $p_1,p_2$.
At these points, tangent vectors are vertical, so $dx = 0$ on all tangent directions.
Thus $\frac{dx}{y}$ appears indeterminate.
Using the relation
\[
    x^2 + y^2 = 1,
\]
differentiate to obtain
\[
    2x\,dx + 2y\,dy = 0,
    \quad
    \Rightarrow \quad
    \frac{dx}{y} = -\frac{dy}{x}.
\]
This expression is regular at $p_1,p_2$.
Hence $\frac{dx}{y}$ is regular everywhere on $X$.

\subsection*{Structure of $\Omega^1[X]$}

We claim
\[
    \Omega^1[X] = \c[X]\frac{dx}{y}.
\]

First, $\frac{dx}{y}$ is regular, and $\Omega^1[X]$ is a $\c[X]$-module, so
\[
    \c[X]\frac{dx}{y} \subseteq \Omega^1[X].
\]

Conversely, let $\omega \in \Omega^1[X]$.
Then $\omega \in \Omega(X)$, and since $\dim \Omega(X) = 1$, we can write
\[
    \omega = f \frac{dx}{y},
    \quad
    f \in \c(X).
\]

Since $\omega$ is regular and $\frac{dx}{y}$ has no poles, we must have $f \in \c[X]$.
Thus
\[
    \Omega^1[X] \subseteq \c[X]\frac{dx}{y}.
\]

Therefore,
\[
    \Omega^1[X] = \c[X]\frac{dx}{y}.
\]